\section{Scope, Method, and Qualifications}

This source-first study tests one reception claim about John 6: whether the
Bread of Life discourse was anciently joined to Bethlehem, the manger, and the
swaddled child. It is not a verse-by-verse commentary on the whole chapter, a
complete history of Eucharistic doctrine, or an adjudication of modern
Johannine composition theories.

English Scripture quotations have been checked at their principal loci against
the repository's registered public-domain Douay--Rheims, whose John 6 numbering
from modern verse 51 onward runs one verse higher. Greek observations are
bounded to the identified Robinson--Pierpont witness and its inline apparatus.
The Vulgate, Pusey's Greek Cyril, the Latin of Augustine and Jerome, the Greek
of Chrysostom, and the manuscript history of pseudo-Gregory have not been
critically collated. Accordingly, the study uses brief paraphrase and makes no
wording-critical conclusion from those witnesses.

The authority classes are distinct: Scripture is the governing text;
patristic witnesses show reception and theological reasoning; later
ecclesiastical teaching may clarify doctrine but is not retrojected into an
earlier author's intention; the final typological synthesis is project
synthesis. The study tests a bounded reception claim rather than claiming a
complete commentary, critical edition, or history of Eucharistic doctrine.

